% Generated draft: requires mathematical and source review.
\documentclass[11pt]{article}
\usepackage[margin=25mm]{geometry}
\usepackage{fontspec}
\setmainfont{Latin Modern Roman}
\usepackage{amsmath,amssymb,mathtools,unicode-math}
\setmathfont{Latin Modern Math}
\usepackage{enumitem,needspace,xurl,hyperref}
\setlength{\parindent}{0pt}
\setlength{\parskip}{5pt}
\setlength{\emergencystretch}{3em}
\newcommand{\N}{\mathbb{N}}
\newcommand{\Z}{\mathbb{Z}}
\newcommand{\Q}{\mathbb{Q}}
\newcommand{\R}{\mathbb{R}}
\newcommand{\C}{\mathbb{C}}
\newcommand{\F}{\mathbb{F}}
\newcommand{\A}{\mathbb{A}}
\newcommand{\PP}{\mathbb{P}}
\newcommand{\E}{\mathbb{E}}
\newcommand{\eps}{\varepsilon}
\newcommand{\dd}{\,\mathrm{d}}
\begin{document}
\section*{UnsolvedMath ID: 30004521}
\textbf{Source problem code:} catalogue entry
\subsection{Definitions and mathematical statement}
For a field \(K\), the projective plane \(\PP^2(K)\) has points given by one-dimensional subspaces of \(K^3\); a \(K\)-defined projective line is the zero locus of a nonzero homogeneous linear form over \(K\). A set \(L\) of \(d\) such lines consists of \(d\) distinct lines. These definitions describe a possible interpretation of the source's terminology; the source does not explicitly specify the allowed fields or the field of definition of the lines.

The supplied problem asks: for every positive integer \(d\), determine the absolute linear Harbourne constant
\[
H(d)=\min_K\;\min_{\lvert L\rvert=d} H(K,L).
\]
Here \(K\) is the source's field parameter and \(L\) is its collection of lines in the projective plane. The displayed expression preserves the supplied formulation, but the record does not define the arrangement invariant \(H(K,L)\), the domain of the outer minimum, or the precise admissibility conditions for the inner minimum. Consequently, this is a statement of the intended computational target with unresolved definitions, rather than a fully specified numerical invariant. The quantifier includes \(d=1\); its treatment must be specified along with the definition of \(H(K,L)\). The use of minima, including their existence and the treatment of empty admissible families, also requires clarification.
\subsection{Short English statement}
Determine the absolute linear Harbourne constant \(H(d)\) for every positive integer \(d\), by minimizing the arrangement invariant \(H(K,L)\) over the intended fields and arrangements of \(d\) projective lines. The supplied record leaves the invariant, admissible fields and arrangements, and boundary conventions unspecified, so a complete precise formulation requires source clarification.
\subsection{Sources}
\begin{itemize}
\item Local UnsolvedMath registry, supplied record: permanent ID \(30004521\), title \emph{Exact Absolute Linear Harbourne Constants}, source problem code \texttt{catalogue entry}, collection \emph{Oberwolfach Reports Open Problems}. The proposed year is recorded as 2020; no individual proposer is supplied.
\item Supplied report citation: \emph{Discrete Geometry} (2021), \emph{Oberwolfach Reports}, DOI 10.4171/owr/2020/30, pp.\ 1469--1528. \url{https://doi.org/10.4171/owr/2020/30}. The background describes Oberwolfach Reports as records of mathematical workshops and their open-problem discussions.
\item Supplied literature reference: Dumnicki--Harrer--Szpond, \emph{On Absolute Linear Harbourne Constants}. \url{https://doi.org/10.1016/j.ffa.2018.03.001}.
\item Supplied literature reference: \emph{Oberwolfach Reports 2020/30}. \url{https://ems.press/content/serial-article-files/46867}.
\item Literature assessment dated exactly \texttt{2026-08-22}: the supplied curation labels the problem open, reports exact values through \(d=31\) and ``for projective-plane orders q squared plus q plus 1,'' and describes the general computation and a conjectural formula as unresolved. This is supplied dated curation derived from the OWR material, not an independently verified literature assessment. The record expressly calls for independent verification of time-sensitive status and cited claims.
\item The registry labels the statement \texttt{corrected\_verified}, but its accompanying verification description says only that formula delimiters, notation, grammar, and typography were repaired while preserving mathematical content and scope. This label does not establish the missing definitions or constitute mathematical verification of this draft.
\end{itemize}
\end{document}
